\chapter{Repository and reproduction}
\label{app:repro}
\chaptermark{Repository and reproduction} %running head

This appendix describes the repository \texttt{fil2405/nesy-nids-bench} and the commands that
rebuild the result files from the data and the tables from the result files. The principles are
those of \cref{sec:harness-reproduce} and \cref{sec:impl-repro}: one run is one configuration and
one seed and writes one result file, and every table is computed from the result files.

\section{Layout}

\Cref{tab:app-layout} lists the directories of the repository and what each holds.

\begin{table}[htbp]
\centering
\footnotesize
\caption{Directories of the repository. The last two are not under version control.}
\label{tab:app-layout}
\setlength{\tabcolsep}{4pt}
\begin{tabular}{@{}p{4.3cm}p{8.9cm}@{}}
\toprule
Directory & Content \\
\midrule
\texttt{src/nesybench/data} & loaders of NSL-KDD, CIC-IDS2017, CIC-DDoS2019 and the transfer schema; split strategies; preprocessing \\
\texttt{src/nesybench/logic} & knowledge base format and degrees of truth, semantics of LTN, edits of E3 and E11, C5 report, post-check \\
\texttt{src/nesybench/models} & rule-only, Random Forest and XGBoost, MLP-CE, 1D-CNN, LTN-sup, LTN-full; batch quotas and epoch selection \\
\texttt{src/nesybench} & runner (\texttt{run.py}), tuning (\texttt{tune.py}), stress configurations (\texttt{stress.py}), statistics (\texttt{compare.py}), metrics, export of models, evaluations of saved models (\texttt{saved\_eval.py}, \texttt{operating.py}, \texttt{latency.py}, \texttt{perturb.py}) \\
\texttt{configs} & one configuration per arm and dataset; \texttt{tuning} the search spaces; \texttt{stress} the configurations of E3, E4, E7 to E9, the transfer, E11 and E13 \\
\texttt{kb} & knowledge bases with their variants, and the effort estimates of E16 \\
\texttt{scripts} & launch scripts of tuning, final, stress and saved-model runs, of the timings of E13 and of the tables \\
\texttt{analysis} & scripts that build tables and figures from the result files \\
\texttt{tests} & unit and integration tests \\
\texttt{thesis} & \LaTeX{} source of this thesis \\
\texttt{data} & downloaded datasets (\texttt{raw}) and cleaned tables (\texttt{interim}) \\
\texttt{results}, \texttt{logs} & result files, saved models, reports and logs \\
\bottomrule
\end{tabular}
\end{table}

\FloatBarrier
\section{Machines}
\label{app:machines}

The runs of the thesis used two machines, named L1 and S1 in the text. Their software is in
\cref{tab:impl-environments}; their build is below.

\subsection{Laptop L1}
\label{app:machine-l1}

L1 is the laptop of the author; \cref{tab:app-machine-l1} gives its build.

\begin{table}[htbp]
\centering
\footnotesize
\caption{Build of L1.}
\label{tab:app-machine-l1}
\begin{tabular}{@{}ll@{}}
\toprule
Model & ASUS TUF Dash F15, FX517ZC \\
CPU & Intel Core i7-12650H \\
GPU & NVIDIA GeForce RTX 3050 Laptop, 4\,GB \\
Memory & 16\,GB, of which 12\,GB assigned to WSL2 \\
Storage & Intel SSDPEKNU512GZ, NVMe, 512\,GB \\
System & Windows 11, Debian under WSL2 \\
\bottomrule
\end{tabular}
\end{table}

\subsection{Server S1}
\label{app:machine-s1}

S1 is a shared machine with one GPU, on which the seeds of a run and the trials of a search ran as
parallel processes (\cref{sec:impl-environment}).
% FIXME(open): provider and components of S1 (CPU, GPU, memory, storage), when the co-supervisor sends them.

\FloatBarrier
\section{Environment and data}

The code needs Python 3.11 or later. The packages are listed in \texttt{requirements.txt}, and
\texttt{requirements.lock} pins the versions of L1 (\cref{tab:impl-environments}).

\begin{lstlisting}[language=bash,numbers=none]
python3 -m venv .venv && source .venv/bin/activate
pip install -r requirements.txt && pip install -e . --no-deps
python scripts/check_env.py
python -m pytest tests
\end{lstlisting}

The datasets go in the folders \texttt{nsl-kdd}, \texttt{cic-ids2017-improved} and
\texttt{cic-ddos2019} of \texttt{data/raw}, or of the folder named by \texttt{NESYBENCH\_DATA}. The
CIC loaders write their cleaned table to \texttt{data/interim} on the first read and reuse it.
Tests that need a dataset are skipped when it is missing.

\section{From the data to the result files}

Each step skips the runs that already have a result file, so a stopped step is resumed by
launching it again. \texttt{[workers]} is the number of runs at the same time.

\begin{lstlisting}[language=bash,numbers=none]
# tuning, then the 10 final seeds with test part and saved model
bash scripts/tune_final.sh <dataset> <arm> [workers]
bash scripts/final_cic.sh <dataset> rule_only
# stress runs E3, E4, E7 to E9, E11, E13 (configurations in configs/stress)
CHECK_ONLY=1 bash scripts/stress_run.sh '<pattern>'
bash scripts/stress_run.sh '<pattern>' [workers]
MODELS_DIR=results/models bash scripts/stress_run.sh 'cic_transfer_*'
# saved models: post-check, E10, E17, operating point, cost curves, latency
bash scripts/saved_runs.sh <dataset> <models>
# E13 timings, one process
bash scripts/e13_time.sh <dataset>
\end{lstlisting}

\texttt{<dataset>} is \texttt{nslkdd}, \texttt{cicids2017} or \texttt{cicddos2019},
\texttt{<models>} the folder of the saved models of the final runs, and
\texttt{<arm>} one of \texttt{mlp\_noweights}, \texttt{mlp\_weights}, \texttt{cnn},
\texttt{ltn\_sup}, \texttt{ltn\_full}, \texttt{rf}, \texttt{xgb}. A pattern selects stress
configurations by name, for instance \texttt{'cicids2017\_ltn\_full\_tuned\_e3\_*'}. The
configurations of the stress runs are written by \texttt{stress.py} from those of the final runs
and are already in the repository.

\section{From the result files to the tables}

One command rebuilds the tables of \cref{ch:results} and of \cref{app:results}:

\begin{lstlisting}[language=bash,numbers=none]
bash scripts/make_tables.sh [results]
\end{lstlisting}

It writes into \texttt{results/tables} one text file per table with the full output, the
\LaTeX{} tables and the figures, and reports the tables whose result files are missing.
\Cref{tab:app-repro-tables} maps the tables of the thesis to its files. The timings of E13 are
printed by \texttt{e13\_time.sh} at its end, the effort estimates of E16 are in
\texttt{kb/effort\_estimate.tsv}, \texttt{analysis/feature\_dictionary.py} writes the tables
of \cref{app:features} and \texttt{analysis/seed\_tables.py} those of every seed in
\cref{app:results}.

\begin{table}[htbp]
\centering
\footnotesize
\caption{Tables of the thesis and the files of \texttt{results/tables} that hold them;
\texttt{<ds>} is the dataset.}
\label{tab:app-repro-tables}
\setlength{\tabcolsep}{4pt}
\begin{tabular}{@{}p{5.4cm}p{7.8cm}@{}}
\toprule
Table & File \\
\midrule
\Cref{tab:res-macro-f1} & \texttt{<ds>\_summary.txt}; 1D-CNN pair in \texttt{<ds>\_cnn.txt} \\
\Cref{tab:res-pairs}, \cref{tab:app-ranks} & \texttt{<ds>\_pairs.txt}, \texttt{<ds>\_pairs.tex}, \texttt{<ds>\_cd.pdf} \\
\Cref{tab:app-epochs} & \texttt{selected\_values.txt}, \texttt{tuning\_epoch\_cap.txt} \\
\Cref{tab:res-rules}, \cref{tab:res-postcheck}, \cref{tab:app-postcheck} & \texttt{<ds>\_postcheck.txt} \\
\Cref{tab:res-sensitivity}, \cref{tab:app-e3}, \cref{tab:app-e4} & \texttt{<ds>\_e3.txt}, \texttt{<ds>\_e4.txt}, \texttt{<ds>\_e11.txt}, \texttt{nslkdd\_e3\_weak.txt} \\
\Cref{tab:res-e7}, \cref{tab:app-e8} & \texttt{<ds>\_e7.txt}, \texttt{<ds>\_e8.txt}, \texttt{nslkdd\_e7\_no\_A3.txt}, \texttt{cicddos2019\_e7\_no\_A2.txt} \\
\Cref{tab:app-e9} & \texttt{cicids2017\_e9.txt} \\
\Cref{tab:res-transfer}, \cref{tab:app-transfer-postcheck} & \texttt{transfer.txt}, \texttt{transfer\_postcheck.txt} \\
\Cref{tab:app-e10}, \cref{tab:app-e17} & \texttt{<ds>\_e10.txt}, \texttt{<ds>\_e17.txt} \\
Inference latency of E12, \cref{tab:app-latency} & \texttt{<ds>\_e12.txt}, \texttt{latency.tex} \\
\Cref{tab:res-operating}, \cref{fig:res-cost-nslkdd} & \texttt{<ds>\_threshold.txt}, \texttt{<ds>\_cost.txt}, \texttt{<ds>\_cost.pdf} \\
\Cref{tab:tuning-selected} & \texttt{selected\_values.txt} \\
Every seed, classes and C5 of the final runs, \cref{tab:app-c5-nslkdd} and the tables beside it & \texttt{back/appendixC\_seeds.tex}, written by \texttt{seed\_tables.py} \\
\bottomrule
\end{tabular}
\end{table}

The result files of the final and stress runs and of the searches, and the saved models of the
final runs, are not in the repository. They are available from the author on request; without them every table needs the runs
of the previous section.

% FIXME(open): the streaming pipeline of the co-supervisor has no directory here yet; add it, its launch script and its commands when it is delivered.
